\section{相关研究}\label{sec:related}
\subsection{图神经网络与图注意力}
GCN将节点特征变换与归一化邻域聚合结合，用于半监督节点分类\cite{kipf2017gcn}。GAT通过邻域内的注意力归一化分配不同邻居的聚合权重，并用多头机制组合表示\cite{velickovic2018gat}。二者说明拓扑与属性可共同进入节点编码。本文采用GraphMAE官方实现中的GAT，不重新选择骨干；GAT内部的注意力权重服务于消息传递，与本文按度数选择训练目标的掩码权重处于不同层次。

\subsection{图自监督的对比式与生成式路线}
图自监督可以通过区分相匹配与受扰动的表示构造目标，也可以通过恢复输入信息进行训练。DGI将局部节点表示与全图摘要联系起来，并通过受扰动图构造对照信号，以互信息最大化思想学习表示\cite{velickovic2019dgi}。这类方法的关键在于构造有效的正负关系；不同方法对增强和训练过程的要求不尽相同，不能将其概括为同一种实现。

生成式方法则明确指定重建对象。例如VGAE使用GCN编码器和内积解码器完成图链接预测\cite{kipf2016vgae}。GraphMAE转向节点特征恢复，并不认为所有图任务都必须首先重建邻接矩阵。重建目标是否适合下游任务、解码器是否能迫使编码器提取有用信息，是两类方法比较时需要区分的因素。

\subsection{掩码自编码与GraphMAE}
BERT通过预测受遮蔽的文本内容开展预训练\cite{devlin2019bert}；视觉MAE使用仅处理可见图像块的编码器和恢复像素的解码器\cite{he2022mae}。二者为从部分观测恢复输入提供了不同领域的例子。GraphMAE沿用掩码重建思想，但适应图数据的邻域依赖：被掩码节点仍进入整图编码，采用GNN解码器，并在解码前重新遮蔽同一批节点的隐表示。其SCE直接衡量特征方向误差，避免将图节点特征机械地视为文本离散词元或图像像素。

本文的度数策略不改变上述信息流，而是改变哪些节点更频繁地成为重建目标。原论文已经指出均匀采样有助于减少邻域观测偏置。因此，偏向高度数节点既有更多上下文可用的动机，也存在破坏上下文与减少低度数监督的风险，需要在统一训练配置下检验。
